\chapter{O papel da energia na circulação de capital}
%\section{A energia enquanto capital circulante}

%\section{Efeitos sobre a rotação do capital}

\hypertarget{os-fundamentos-da-maquinaria}{%
\section{Os fundamentos da maquinaria}\label{os-fundamentos-da-maquinaria}}

N'O Capital, Marx descreve a composição de ``toda maquinaria desenvolvida'' como consistindo de 1) uma ``\textit{máquina-motriz}'' --- que ``atua como força motora do mecanismo inteiro'' ---, 2) um ``\textit{mecanismo de transmissão}'', o qual ``regula o movimento, modifica sua forma onde é necessário (\ldots) e o distribui e transmite'' à sua mais importante parte, 3) a ``\textit{máquina-ferramenta}'', por meio da qual uma máquina ``se apodera do objeto de trabalho e o modifica conforme a uma finalidade'' \parencite[, p. 446-7, grifo meu]{Marx2017}.

Tais metálicos meios de produção somente tornam-se \textit{máquinas} propriamente ditas quando não são mais meros meios de trabalho, e sim quando assumem o papel primário que antes era ocupado pelo trabalhador no processo de produção, restringindo-se o fator subjetivo do trabalho agora a ``vigiar a máquina com os olhos e corrigir os erros dela com as mãos'' \parencite[, p. 448]{Marx2017}. Isso se dá pelo desenvolvimento da divisão do trabalho e da especialização não só dos trabalhos particulares, como também, e especialmente, das ferramentas empregadas nos trabalhos particulares. Este é o fundamento da mecanização do trabalho: ``Quando cada processo é reduzido ao uso de uma ferramenta simples'' (em decorrência do avanço da divisão do trabalho), ``a união de todas essas ferramentas, acionadas por uma força motriz, constitui uma máquina'' \parencite[, p. 136 (§167)]{Babbage2009}.

O processo de ``maquinaria e grande indústria'', porém, não consiste somente na mecanização de processos parciais de trabalho, e sim da automatização\footnote{Ou melhor: do máximo possível que for automatizável.} de processos inteiros de produção. A divisão social do trabalho entre trabalhadores parciais, ``a cooperação peculiar à manufatura'', aparece agora como uma divisão do trabalho entre máquinas individuais; as ferramentas especializadas dos trabalhadores parciais aparecem agora como ``ferramentas de máquinas de trabalho especializadas''; o processo de produção, que antes girava em torno do trabalhador, aparece agora como um sistema de produção que gira em torno da maquinaria \parencite[, p. 453]{Marx2017}. Em particular, quando o processo consegue executar ``todos os movimentos necessários ao processamento da matéria-prima sem precisar da ajuda do homem, mas apenas de sua assistência'', este processo, segundo Marx, descreve um ``sistema automático de maquinaria'' \parencite[, p. 455]{Marx2017}.

Dessa forma, o surgimento da máquina destaca o trabalho subsumido pelo capital, ao mesmo tempo que o ofusca. Destaca o trabalho subsumido pelo capital pois, assim como o capital emprega mãos individuais mas usufrui do trabalhador coletivo\footnote{``O que caracteriza a divisão do trabalho? Que o trabalhador parcial \textbf{não} produz mercadoria. Apenas o \textit{produto comum} {[}``\textit{gemeinsame Produkt}''{]} dos trabalhadores parciais converte-se em mercadoria.'' \parencite[, p. 429, grifo e adendo meus]{Marx2017}.}, o capital compra partes individuais de máquinas, mas usufrui do produto da linha inteira de montagem. Ofusca o trabalho subsumido pelo capital pois, assim como o trabalho pretérito ofusca o trabalho vivo empregado, o sistema automatizado de produção ofusca a divisão social do trabalho à qual este sistema remete. As funcionalidades da máquina individual, seu \textit{modus operandi}, seu \textit{output}, seu encaixe no resto da ``fábrica''', tudo isso torna-se sua natural \textit{raison d'être}; torna-se seu papel a desempenhar dentro do enorme ecossistema industrial, abstraindo-se de todo o processo técnico-social através do qual se deu sua ``seleção social'' dentre várias possíveis alternativas no mercado. É cativante, portanto, ver o aspirante a capitalista deixar-se levar pelo \textit{curriculum vitae} da máquina que ele deseja empregar, ao mesmo tempo que se esquece que, assim como as tartarugas de Galápagos, máquinas não conseguem se levantar sozinhas quando tombam, \textit{crasham} ou pifam --- por \textit{algum} motivo desconhecido e diabólico, a Natureza faz com que ele precise ter técnicos especializados sempre a postos para cada autômato que lhe brilhe aos olhos.\footnote{E o mesmo vale, hoje em dia, para o algoritmo e a tão falada \textit{IA} --- ``o que quer que isso seja'', murmura o capitalista para si próprio, enquanto conta as cédulas do dinheiro que reavê após suas demissões em massa para a compra de \textit{tokens} de sua \textit{hyperscaler} de escolha. Mesmo aqui, desafortunadamente, sempre chega o momento em que \textit{bugs} têm de ser resolvidos por empregados de dentro da corporação, invés de por serviços contratados de fora.}

\hypertarget{o-contexto-de-surgimento-da-energia-abstrata-mercadoria-energia}{%
\section{O contexto de ``surgimento'' da energia abstrata (mercadoria-energia)}\label{o-contexto-de-surgimento-da-energia-abstrata-mercadoria-energia}}

\textcite{Pasquinelli2023} elabora que \textgreater{} ``Não foi a invenção da máquina a vapor (\textit{meio de produção}) que desencadeou a Revolução Industrial (como é comum teorizar no discurso ecológico), mas sim o desenvolvimento do capital e do trabalho (\textit{relações de produção}) que exigiam uma fonte de energia mais potente.'' \parencite[, p. 107]{Pasquinelli2023}.

Neste contexto, ele argumenta, seguindo \textcite{Malm2016}, que \textgreater{} ``o carvão contribuiu para a aceleração do capitalismo industrial não apenas por seu potencial energético, mas porque suas propriedades físicas, como leveza, homogeneidade e mensurabilidade, correspondiam perfeitamente às novas dimensões abstratas do capital. As máquinas a vapor substituíram os moinhos de água \textit{não porque o carvão fosse mais barato e abundante que a água}, mas porque fornecia um \textit{fluxo de energia mais estável} do que a chuva e permitia que as fábricas se instalassem \textit{perto das áreas urbanas}, onde a classe trabalhadora vivia na época.'' \parencite[, p. 122-3, grifo meu]{Pasquinelli2023}

Em suma, o carvão suplantou as rodas d'água das indústrias inglesas não por ter sido, desde sempre, uma fonte energética mais eficiente --- de fato, não o foi até meados do século XIX\footnote{``O vapor era mais barato que a água? À primeira vista, essa é a explicação mais plausível: os capitalistas do algodão optaram pelo vapor porque um cavalo-vapor dele era mais barato que um de água. Uma roda d'água representava um investimento considerável. A própria roda precisava ser comprada, instalada em uma casa de máquinas {[}\textit{wheel-house}{]} e, na maioria dos casos, complementada por uma represa para garantir um abastecimento regular de água. Em seguida, o dono da fábrica precisava construir um sistema de condutos --- canais, valas, comportas --- para levar a água até a roda na quantidade adequada, presumindo que fosse do tipo padrão de queda superior ou de peito {[}\textit{standard overshot or breast type}{]}. Uma máquina a vapor, por outro lado, era composta de ferro, latão e cobre, volante {[}\textit{flywheel}{]}, caldeira e tubulações; fixada em uma estrutura sólida em uma casa de máquinas {[}\textit{engine-house}{]} especial, sua construção exigia mão de obra especializada. Além disso, havia a necessidade ocasional de reparos extensivos após avarias e as taxas de depreciação extremamente altas, enquanto as rodas d'água podiam funcionar com apenas manutenção mínima por décadas ou até mesmo um século. A água jorrava de graça. Uma vez que o capitalista garantisse um arrendamento do proprietário da terra, pagando aluguel pelo direito de utilizar o riacho, não haveria mais custos com combustível. O carvão precisava ser comprado constantemente no mercado. A soma dessas relações é amplamente aceita na literatura: \textit{as rodas d'água eram consistentemente mais baratas por cavalo-vapor do que as máquinas a vapor} {[}grifo no original{]} no início do século XIX. `É difícil resistir à conclusão', escreve o historiador do algodão Stanley Chapman, `de que o vapor era mais caro do que as instalações hidráulicas mais dispendiosas.'\,'' \parencite[, p. 29-30]{Malm2013}.}. A vantagem dos motores a vapor não eram sua eficiência energética imediatamente considerada, e sim a viabilidade que eles traziam de estabelecer indústrias próximas a centros urbanos; o motor a vapor, conforme Malm, ``não abriu novas reservas de energia tão necessárias, mas sim deu acesso a mão de obra explorável'' \parencite[, p. 33]{Malm2013}.

Jevons, em \textit{The Coal Question}, também reflete sobre essa questão: \textgreater{} ``A necessidade, mais uma vez, de levar o trabalho até a usina, e não a usina até o trabalho, é uma desvantagem na energia hidráulica e impede completamente a concentração do trabalho em uma mesma região, que é {[}esta concentração{]} altamente vantajosa para o aperfeiçoamento do nosso sistema mecânico.'' \parencite[, p. 89]{Jevons1866}

Pasquinelli, portanto, resume que ``se o carvão passou a ser utilizado em todo o espectro da produção'' (vide acima), ``foi porque era a fonte mais adequada de \textit{energia abstrata}, onde `abstrato' significa facilmente calculável em termos de custo, transporte, estoque, desempenho e organização social'' \parencite[, p. 122-3, grifo no original]{Pasquinelli2023}. O carvão tornou-se, contudo, mais do que apenas energia abstrata: tornou-se, de fato, uma mercadoria com o valor de uso específico de \textit{ser uma fonte de força-motriz} a máquinas-ferramentas\footnote{Evidentemente, força-motriz somente para máquinas-ferramentas que ``comportem'' seu usufruto, i.e.~que possuam ``mecanismos de transmissão'' adequados. Máquinas que estivessem fixadas em usinas movidas a rodas d'água não podiam meramente ``passar a usar'' motores a vapor, pois isso requer mobilização de capital fixo novo.}, pois, ao contrário da água, que era uma fonte de força-motriz que ``vinha fluindo gratuitamente'', o carvão tinha de ser ``constantemente comprado no mercado'' \parencite[, p. 29]{Malm2013}; o carvão, ao contrário da água, não era dado gratuitamente pela Mãe Natureza, sendo fruto necessariamente do suor e do sofrimento.\footnote{``E ao homem {[}Deus{]} declarou: `{[}\ldots{]} com sofrimento você se alimentará {[}da terra{]} todos os dias da sua vida. Ela lhe dará espinhos e ervas daninhas, e você terá que se alimentar das plantas do campo. Com o suor do seu rosto, você comerá o seu pão, até que volte à terra, visto que foi dela tirado; porque você é pó e ao pó voltará.'\,'' (Gn, 3:17-19).}

\hypertarget{os-impactos-da-maquinariamercadoria-energia-sobre-o-processo-de-produuxe7uxe3o-do-capital}{%
\section{Os impactos da maquinaria/mercadoria-energia sobre o processo de produção do capital}\label{os-impactos-da-maquinariamercadoria-energia-sobre-o-processo-de-produuxe7uxe3o-do-capital}}
Compreender o impacto que novas tecnologias desempenham sobre processos de geração de valor (i.e. processos de produção) requer tomar seus possíveis impactos \textit{in abstractu}, a fim de melhor aferir suas consequências isoladamente para, depois, verificar seu impacto em conjunto.

\hypertarget{produtividade-do-trabalho}{%
\subsection{Produtividade do Trabalho}\label{produtividade-do-trabalho}}

\begin{quote}
``{[}A{]} produtividade sempre diz respeito ao trabalho. O nível de produtividade nos informa uma determinada relação entre o trabalho vivo executado e o produto.''

``O ganho de produtividade refere-se à diminuição do trabalho vivo exigido para produzir-se uma mercadoria (ou qualquer quantidade determinada de mercadorias). Como estão abstraídas {[}a princípio{]} variações de eficiência, o conjunto de matérias-primas e matérias auxiliares consumidos continua sendo o mesmo para qualquer nível dado de produção.'' \parencite[p. 233]{SaBarreto2021a}
\end{quote}

O aumento da produtividade do trabalho permite com que o dispêndio de força de trabalho para a produção de uma mercadoria seja reduzido. Isso é potencializado com a proliferação da maquinaria.

Suponhamos a produção de \(1\) unidade de uma dada mercadoria \(M\), em que o mercado produz, em média, em \(\Delta t\), e um capital mais produtivo a produz em \(\Delta t^{\prime}<\Delta t\). Demais circunstâncias são as mesmas: ambos despendem \(c\) nesta produção unitária, à mesma composição orgânica do capital $c/v$. Cabe analisarmos, neste caso, somente os dispêndios distintos de tempo. O mais-valor extra que este capital mais produtivo afere tem magnitude de tempo \(\Delta t-\Delta t^{\prime}\); ou seja, o mais-valor total que ele afere é seu mais-valor usual --- dado pelo tempo excedente de trabalho extraído na jornada de trabalho de seu processo de produção ---, mais este ``tempo extra''.

Como um exemplo mais geral, tenhamos agora que a produção de \(N\) unidades de uma dada mercadoria \(M\) leve, em média, \(\Delta t\) e requeira adiantamento de capital \(c+v\), e que seja produzida (e vendida) a \(c+(1+m^{\prime})v\),\footnote{O valor social também é constituído pelo mais-valor \(m^{\prime}v\), pois ele é extraído \textit{dos trabalhadores, pelo capitalista}! Neste setor, \textit{todos} os capitalistas extraem este mais-valor --- de novo, em média --- \textit{de seus respectivos empregados assalariados}.} onde a taxa de mais-valor \(m^\prime\) é dada para todos; suponha-se também que a composição orgânica do capital \(\frac{c}{v}\) seja a mesma para todos os capitais deste setor; deixe-se explícito que este dispêndio de capital variável \(v\) diz respeito \textit{a este período médio de tempo \(\Delta t\)}. Tenhamos agora algum capital que tenha produtividade acima da média, portanto produzindo, digamos, \(N+\Delta N\) unidades neste mesmo tempo \(\Delta t\). \textbf{Na média}, o valor unitário da mercadoria \(M\) é \(\frac{c+(1+m^{\prime})v}{N}\), enquanto para este capital mais produtivo é \(\frac{(1+\sigma)c+(1+m^{\prime})v}{N+\Delta N}\), pois 1) a produção de \(N\) mercadorias requer, como na média, o mesmo adiantamento de capital \(c+v\), porém é produzida em \(\Delta t^{\prime}<\Delta t\); 2) a produção de \(\Delta N\) mercadorias, no período \(\Delta t-\Delta t^{\prime}>0\), requer um adiantamento a mais de capital constante --- dado por essa proporção $\sigma > 0$ ---, \textit{porém emprega o mesmo capital variável adiantado \(v\)}, pois ele foi contratado pelo período de tempo \(\Delta t\) \textit{inteiro}, não somente \(\Delta t^{\prime}\). Portanto, no período de tempo \(\Delta t\) inteiro, o capital adiantado para este capital mais produtivo será \((1+\sigma)c + v\), produzindo \(N+\Delta N\) mercadorias, cujo valor total produzido será \((1+\sigma)c + (1+m^{\prime})v\).

Para que compense para este capital produzir a este nível maior de produtividade, é necessário que seu ``custo'' (valor ``individual'') seja menor (ou, no máximo, igual) que o ``valor social''\footnote{A rigor, ``valor social'' é redundante: o valor de uma mercadoria é definido pelo tempo de trabalho \textit{socialmente} necessário para sua produção. Neste caso, o que se chama de ``valor individual'' trata-se de valor 1) que foi adiantado pelo capitalista em capital constante $c$, o qual é \textit{transferido} à mercadoria final, e 2) que foi \textit{produzido} pelo capital variável em ação, o qual reproduz o seu próprio valor $v$, que foi adiantado pelo capitalista, mais um valor excedente $m$ (mais-valor) produzido no resto da jornada de trabalho em que esta força de trabalho foi empregada.}: 
\begin{align*}
(\cancel{ 1 }+\sigma)c + \cancel{ (1 +m^{\prime})v } &\leq \cancel{ c }+\cancel{ (1+m^{\prime})v } + \frac{\Delta N}{N}(c+(1+m^{\prime})v) \\
\iff \sigma c & \leq \frac{\Delta N}{N}[c+(1+m^{\prime})v]
\end{align*}

Dito em termos ``vulgares'': o ``custo extra'' (\(\sigma c\)) tem que ser menor (ou igual) ao ``lucro extra'' \(\frac{\Delta N}{N}(c+v)\) de \(\Delta N\) mercadorias a mais sendo vendidas ao valor social unitário, \(\frac{c+v}{N}\), para que o capitalista possa aderir (racionalmente, claro) a esta ``estratégia'' de produção. 

\hypertarget{intensidade-do-trabalho}{%
\subsection{Intensidade do Trabalho}\label{intensidade-do-trabalho}}
\textbf{Muito} mais complicado do que parece. Marx o coloca como um ``grau de condensação'' da jornada de trabalho, em que ``A hora mais intensa da jornada de trabalho de 10 horas encerra tanto ou mais trabalho, isto é, força de trabalho despendida, que a hora mais porosa da jornada de trabalho de 12 horas.'' \parencite[p. 482-3]{Marx2017}. Os cálculos, porém, são muito nebulosos e, na literatura marxista, basicamente inexistentes. Obtê-los (quase que) \textit{ex nihilo} requer tempo, mas estou quase lá!


\hypertarget{eficiuxeancia}{%
\subsection{Eficiência}\label{eficiuxeancia}}

\begin{quote}
``Eficiência sempre estará relacionada aos meios de produção. O nível de eficiência nos informa uma determinada relação entre insumos e produto. Quanto mais próxima essa razão estiver de \(0\), mais eficiente é o processo.'' \parencite[, p. 60]{SaBarreto2022}
\end{quote}

Um capital possui uma eficiência maior do que a média quanto requer menos capital constante para produzir uma mesma quantidade \(N\) de uma dada mercadoria \(M\) em dado tempo \(\Delta t\). Se, por exemplo, em média adianta-se \(c+v\) neste período e produz-se \(c+(1+m^{\prime})v\), então um capital será mais eficiente se produzir a mesma quantidade neste período, mas precisar adiantar somente \((1-\sigma)c + v\), com \(\sigma>0\). Como produz \((1-\sigma)c + (1+m^{\prime})v\) e vende por \(c+(1+m^{\prime})v\), este capital consegue extrair mais-valor extra \(\sigma c\).\footnote{Como capital constante somente \textit{transfere} seu valor ao produto, fica claro aqui que \(\sigma c\) é justamente a parcela de valor que deixa de ser transferida pelo capital mais eficiente, \textit{mas que é transmitida pelo capital médio}. Já aqui é possível entrever o papel da transferência de valor dos capitais menos ``eficientes'' aos mais ``eficientes''; não é claro a este nível, porém, que é a composição orgânica do capital \(\frac{c}{v}\) que dita a direção da ``sifonagem'' destes diferenciais de valor. Tal nível de abstração somente é elaborado por Marx no livro III \parencite{Marx2017a}.}

Olhando a situação em um nível mais concreto, o capitalista enxerga seu capital não como composto por fatores subjetivos e objetivos, por capital variável e capital constante, e sim como \textit{insumos}, como fatores de produção; sua composição, que aparece como capital fixo e capital circulante, diferencia-se somente pela ``perenidade'' de seus insumos ao longo de rotações do capital\footnote{Não à toa esta diferenciação deve muito aos fisiocratas, que tiraram da produção agrícola de culturas \textit{perenes} e culturas \textit{anuais} a diferenciação de capital \textit{fixo} e capital \textit{circulante}, respectivamente. \#to-be-elaborated checar anotação de \textcite{Marx2014} ou \textcite{Marx2017a} sobre esta menção dos fisiocratas.}, e não de seu papel na produção ou transferência de valor. Não é à toa, portanto, que a economia tradicional utiliza as categorias de produtividade e eficiência intercambiavelmente: no que diz respeito a seu capital adiantado, ``usar menos'' insumos metálicos ou humanos não faz diferença para o capitalista em ação. Ele só consegue entrever a diferença destas categorias justamente quando \textit{não} compensa adotar um plano ``tecnicamente mais eficiente'' de produção, em cujo caso a adoção desta tentadora ``inovação tecnológica'' tem de ser resignadamente rechaçada pela mesma estoica abstinência com que acredita que se abstém dos prazeres da carne em prol do \textit{business as usual}.